% Blueprint content for the prime_gap formalisation.
%
% Conventions used throughout this file:
%   \rocq{PrimeGapS1.<File>.<name>}  -- link to a Rocq declaration
%   \rocqok                          -- the surrounding env is fully Qed-formalised
%   \uses{label1, ...}               -- dep-graph dependencies
%   \label{...}                      -- unique label, prefixed by env type
%
% The Rocq logical name of every theory file is `PrimeGapS1.<File>` (see
% `_CoqProject`), so e.g. `CertRayleigh.maynard_M105_certified_rayleigh` becomes
% `PrimeGapS1.CertRayleigh.maynard_M105_certified_rayleigh`.

%%%%%%%%%%%%%%%%%%%%%%%%%%%%%%%%%%%%%%%%%%%%%%%%%%%%%%%%%%%%%%%%%%%%%%%%
\chapter*{Introduction}

This blueprint accompanies a Rocq formalisation that machine-verifies the
numerical step underpinning Maynard's bounded-gap theorem:
\[
   \Mk[105] \;=\; 105 \cdot \sup_F  J_{105}(F) / I_{105}(F) \;>\; 4.
\]
The headline is formula~(8.15) of James Maynard,
\emph{Small gaps between primes} (arXiv:1311.4600;
Annals of Mathematics \textbf{181} (2015), 383--413).
In the published proof, this inequality is checked by a Mathematica
notebook (\texttt{Computations.nb}, supplied as supplementary material
with the arXiv preprint). The present project replaces that
Mathematica-based computation with a Rocq proof that re-derives every
algebraic claim inside the kernel.

The proof strategy is the same as Maynard's notebook: instead of
computing $\lammax(\Mone^{-1}\Mtwo)$ inside the kernel, we exhibit
a single 42-entry rational \emph{witness vector}
$v_{\mathrm{witness}} \in \rat^{42}$ (shipped in
\texttt{Witness\_Quad.v}, obtained by snapping the top eigenvector
of $\Mone^{-1}\Mtwo$ to small-denominator rationals via
continued-fraction convergents) and verify the strict Rayleigh-
quotient bound
\[
4 \cdot v^T \Mone\, v   \;<\;   105 \cdot v^T \Mtwo\, v
\]
at $v = v_{\mathrm{witness}}$ in pure integer arithmetic by a single
\vmcompute{} reflexivity. By Maynard's Lemma~8.3, this entails
$\Mk[105] > 4$ (since the supremum over admissible test functions is
at least the quotient at any individual choice).

The end-to-end Rocq statement is
\begin{verbatim}
Theorem maynard_M105_certified_rayleigh :
  (forall i j : nat, (i < 42)%nat -> (j < 42)%nat ->
     M1_spec_ij i j = Z2rat (mat_get M1_int i j) / Z2rat D_M1) /\
  (forall i j : nat, (i < 42)%nat -> (j < 42)%nat ->
     M2_spec_ij i j = Z2rat (mat_get M2_int i j) / Z2rat D_M2) /\
  4%:Q * quad_spec M1_spec_ij < 105%:Q * quad_spec M2_spec_ij.
\end{verbatim}
proven in \texttt{theories/S1/CertRayleigh.v}, where
$\texttt{Z2rat}\,(z : \Z) : \rat := (\texttt{Z\_to\_int}\,z)\%\texttt{:\~{}R}$
is the thin embedding of $\Z$ into $\rat$. The first two conjuncts
are each a single composed identity per matrix: the $\rat$-level
paper-form spec $\texttt{M\{1,2\}\_spec\_ij}$ (the readable
transcription of Maynard's Lemma~8.2) is equal, as a rational, to
the FLINT-shipped integer entry
$\texttt{mat\_get}\,M_{1,2}^{\mathrm{int}}\,i\,j$ divided by the
common denominator $\DMone$ / $\DMtwo$. The third conjunct is the
strict Rayleigh-quotient bound at the shipped witness, expanded as
a bigop on the paper-form spec matrices.

\paragraph{Scope.}
What this Rocq layer proves, kernel-checked, is the left-hand side of
\[
\underbrace{\;v^T \Mone v < (105/4) \cdot v^T \Mtwo v \text{ at some } v \in \rat^{42}\;}_{\text{Rocq, kernel-checked}}
\quad\Longrightarrow\quad
\underbrace{\;\Mk[105] = 105 \cdot \sup_F J_{105}(F)/I_{105}(F) \;>\; 4\;}_{\text{Maynard's paper, Lemma 8.3}}.
\]
The right arrow is Maynard's Lemma~8.3 (a generalised Rayleigh-quotient
identity proved in the Annals paper, refereed there); the project does
not re-formalise this analytic step. The project's scope is exactly to
replace the Mathematica computation, not to re-formalise Maynard's paper.

\paragraph{Status.}
The headline theorem and every intermediate lemma is \Qed{} and
\texttt{Print Assumptions} reports \emph{Closed under the global
context}. The proof chain is axiom-free end-to-end: no
\texttt{Admitted}, no \texttt{Axiom}, no \texttt{Parameter}
anywhere in \texttt{theories/S1/}. Because every reduction is in
$\Z$ arithmetic (\vmcompute{} on \texttt{list (list Z)} matrices
and \texttt{list (Z * Z)} vectors), the headline does not even pull
in the native 63-bit primitive-integer interface: no
\texttt{PrimInt63}, no \texttt{Uint63Axioms}, no \texttt{CarryType}
appear in the reported assumptions.

\paragraph{Two certified routes.}
The repository ships the conclusion $\Mk[105] > 4$ in two axiom-free
forms. \emph{Route~1} (Chapters~\ref{ch:matrix}--\ref{ch:rayleigh}) is the
direct Rayleigh-quotient witness described above, with headline
\rocq{PrimeGapS1.CertRayleigh.maynard_M105_certified_rayleigh}.
\emph{Route~2} (Chapter~\ref{ch:eigen}) restates the result in Maynard's
own spectral language --- $M_{105} = 105 \cdot \Mone^{-1}\Mtwo$ has an
eigenvalue $> 4$ over the algebraic complex numbers --- with headline
\rocq{PrimeGapS1.MaynardEigen.maynard_M105_certified}; it is built on
Route~1's inequality and additionally constructs an integer characteristic
polynomial, a Chinese-remainder lift, and an explicit eigenvalue, all over
$\Z$ and $\Fp$. Both are \Qed{} and report \emph{Closed under the global
context}.

\paragraph{How to read this blueprint.}
Each chapter mirrors a Rocq source file or a small group of related
files. Inside each chapter, mathematical statements are presented as
\texttt{lemma} / \texttt{theorem} / \texttt{definition} environments,
each tagged with the corresponding Rocq declaration via
\verb|\rocq{...}|. The \verb|\rocqok| marker means the surrounding
environment is fully formalised and closed by \verb|Qed| in the source
tree. \verb|\uses{...}| declarations populate the dependency graph
shown in the web version of the blueprint.

%%%%%%%%%%%%%%%%%%%%%%%%%%%%%%%%%%%%%%%%%%%%%%%%%%%%%%%%%%%%%%%%%%%%%%%%
\chapter{The matrix construction}\label{ch:matrix}

This chapter sets up the input data that the rest of the proof consumes:
the 42-element basis, the closed-form entries of $\Mone$ and $\Mtwo$,
the $\Z$-level cross-multiplication checks that anchor the
FLINT-shipped integer matrices to those closed forms, the kernel-Qed
bridge between the rat-level paper-form spec and the Z-level
computational spec, and the bridge into MathComp's $\Q$-level matrices.
All material here lives inside the files \texttt{MaynardSpec.v},
\texttt{MaynardBasis.v}, \texttt{MaynardFactQ.v}, the
\texttt{MaynardVerify/} directory (\texttt{Def.v}, six parallel
\texttt{M2\_0..5.v} chunks, and the assembly \texttt{MaynardVerify.v}),
and \texttt{MaynardSpecBridge.v}. Throughout this chapter we follow the
v3/Annals numbering of Maynard's paper: Lemma~8.1 / 8.2 / 8.3 (=
Lemma~7.1 / 7.2 / 7.3 in v1).

\section{Conventions}

\begin{convention}
The two power sums are
$P_1 = \sum_i t_i$ and $P_2 = \sum_i t_i^2$. The basis monomials are
$F_{b,c} = (1 - P_1)^b \cdot P_2^c$ with $b + 2c \le 11$. Indices
$(b_i, c_i)$, $(b_j, c_j)$ denote bidegrees of the $i$-th and $j$-th
basis monomials. The outer integration dimension is $k = 105$; we use
two constants $K_1 = 105$ (outer simplex) and $K_2 = 104 = K_1 - 1$
(residual simplex after the eq.~8.8 substitution that integrates out
$t_1$).
\end{convention}

\begin{definition}[Constants $K_1$, $K_2$]\label{def:K1K2}
\rocq{PrimeGapS1.MaynardSpec.K1, PrimeGapS1.MaynardSpec.K2}
\rocqok
$K_1 := 105$, $K_2 := 104$. The first is the outer simplex dimension
(used in $M_1$ and in the threshold $4/K_1$); the second is the
residual simplex dimension after the eq.~8.8 substitution (used in $M_2$).
\end{definition}

\section{The 42-element basis}

\begin{definition}[Basis]\label{def:maynard_basis}
\rocq{PrimeGapS1.MaynardBasis.maynard_basis}
\rocqok
The Maynard basis is the explicit list of 42 pairs $(b, c) \in \N^2$
with $b + 2c \le 11$, in the Mathematica enumeration order used by
\texttt{flint\_probe.py}:
\[
[(0,0), (1,0), (0,1), (2,0), (1,1), (3,0),\, (0,2), \ldots, (9,1), (11,0)].
\]
\end{definition}

\begin{lemma}[Basis size]\label{lem:basis_size}
\rocq{PrimeGapS1.MaynardBasis.maynard_basis_size}
\rocqok
\uses{def:maynard_basis}
$\#\{(b, c) \in \N^2 : b + 2c \le 11\} = 42$.
\end{lemma}
\begin{proof}\rocqok
$\sum_{c=0}^{5}(12 - 2c) = 12 + 10 + 8 + 6 + 4 + 2 = 42$.
A single \vmcompute{} on the list closes it.
\end{proof}

\begin{lemma}[Basis matches FLINT enumeration]\label{lem:basis_matches}
\rocq{PrimeGapS1.MaynardBasis.maynard_basis_eq_witness}
\rocqok
\uses{def:maynard_basis}
$\texttt{maynard\_basis} = \texttt{Witness.basis}$,
i.e.\ the hand-readable basis in \texttt{MaynardBasis.v} agrees
literally with the autogenerated basis used to index the shipped
integer matrices in \texttt{Witness.v}.
\end{lemma}
\begin{proof}\rocqok
\vmcompute{} on a list of 42 pairs.
\end{proof}

\begin{lemma}[Basis is the canonical set]\label{lem:basis_spec}
\rocq{PrimeGapS1.MaynardBasis.maynard_basis_spec}
\rocqok
\uses{def:maynard_basis}
For every $p \in \N \times \N$,
\[
  p \in \texttt{maynard\_basis} \quad\iff\quad p.1 + 2 \cdot p.2 \le 11.
\]
\end{lemma}
\begin{proof}\rocqok
The basis is permutation-equal to the canonical filter
\[
  \texttt{canonical\_basis} \;:=\; \bigl[\, p \,\big|\, p \in
    \texttt{iota}\,0\,12 \times \texttt{iota}\,0\,6,\; p.1 + 2\,p.2 \le 11
  \bigr]
\]
by a \vmcompute{} \texttt{perm\_eq} lemma, and \texttt{canonical\_basis}
satisfies the spec by construction (\texttt{mem\_filter} +
\texttt{allpairs\_f}; the residual $2c \le 11 \Rightarrow c < 6$ step
closes by \texttt{lia}).
\end{proof}

\begin{lemma}[Basis has no duplicates]\label{lem:basis_uniq}
\rocq{PrimeGapS1.MaynardBasis.maynard_basis_uniq}
\rocqok
\uses{def:maynard_basis}
$\texttt{uniq}\;\texttt{maynard\_basis}$.
\end{lemma}
\begin{proof}\rocqok
\vmcompute{} on the 42-pair list.
\end{proof}

\begin{remark}
Lemmas~\ref{lem:basis_size}, \ref{lem:basis_spec},
and~\ref{lem:basis_uniq} together pin \texttt{maynard\_basis} to
exactly the multiset $\{(b, c) \in \N^2 : b + 2c \le 11\}$. A reviewer
never has to inspect the literal 42-pair list — the basis is fully
characterised by the predicate $b + 2c \le 11$, the size $42$, and the
absence of duplicates. Order matters because the matrix rows and
columns and the entries of the rational witness vector
$v_{\mathrm{witness}}$ (Definition~\ref{def:v_witness}) must all
align with the same indexing in $M_{1,2}^{\mathrm{int}}$.
\end{remark}

\begin{remark}
This is a lower-bound subspace argument. Maynard's full optimisation is
over all symmetric $F$ of degree $\le 11$ in $k = 105$ variables;
restricting to polynomials in the two power sums $P_1, P_2$ is a
strict subspace, and the supremum of a Rayleigh quotient over a
subspace is $\le$ the unrestricted supremum. So a lower bound of
$4/105$ on the 42-dimensional restricted problem yields
$\Mk[105] > 4$.
\end{remark}

\section{The polynomial \texorpdfstring{$G_{n,2}(k)$}{G\_n,2(k)} (Lemma 8.1)}

\begin{definition}[Compositions of $n$]\label{def:compositions}
\rocq{PrimeGapS1.MaynardSpec.compositions}
\rocqok
For $r, n \in \N$,
$\Comp(r, n) := \{(b_1, \ldots, b_r) \in \N_{\ge 1}^r : \sum_s b_s = n\}$,
returned as a list in lexicographic order.  For $r = 0$ this is the
singleton $\{[\,]\}$ if $n = 0$ and empty otherwise; for $r = n + 1$
or larger it is empty.
\end{definition}

\begin{definition}[Inner factor $\Cff(a)$]\label{def:cff}
\rocq{PrimeGapS1.MaynardSpec.cff}
\rocqok
\uses{def:compositions}
For a composition $a = (b_1, \ldots, b_r)$,
\[
\Cff(a) \;:=\; \prod_{s=1}^{r} \frac{(2 b_s)!}{b_s!}.
\]
This is exactly the per-composition inner product appearing in
Maynard's Lemma~8.1 (with $j = 2$, so $j b_s = 2 b_s$).
\end{definition}

\begin{definition}[Maynard's $G_{n,2}(k)$]\label{def:G_2}
\rocq{PrimeGapS1.MaynardSpec.G_2}
\rocqok
\uses{def:compositions, def:cff}
For $n, k \in \N$:
$\Gnk{0}{k} = 1$, and for $n \ge 1$,
\[
\Gnk{n}{k} \;=\; n! \cdot \sum_{r = 1}^{n} \binom{k}{r}
                          \sum_{a \in \Comp(r, n)} \Cff(a).
\]
This is Lemma~8.1 of arXiv:1311.4600~v3 character-for-character: the
$n!$ prefactor, the outer $r$-sum, the binomial coefficient, the
inner sum over length-$r$ compositions of $n$, and the inner product
$\prod_s (j b_s)!/b_s!$ with $j = 2$ all appear at the same positions.
\end{definition}

\begin{lemma}[Sanity values of $G_{n,2}$]\label{lem:G_2_small}
\rocqok
\uses{def:G_2}
$\Gnk{0}{k} = 1$, $\Gnk{1}{k} = 2k$, $\Gnk{2}{k} = 4 k^2 + 20 k$.
\end{lemma}
\begin{proof}\rocqok
$\Gnk{1}{k}$: the only length-1 composition of $1$ is $[1]$, with
$\Cff([1]) = 2!/1! = 2$, so the sum is $1! \cdot \binom{k}{1} \cdot 2 = 2k$.
$\Gnk{2}{k}$: length-1 contributes $\binom{k}{1} \cdot \Cff([2]) = k \cdot 4!/2! = 12k$;
length-2 contributes $\binom{k}{2} \cdot \Cff([1,1]) = \binom{k}{2} \cdot (2 \cdot 2) = 2k(k-1)$;
total $= 2! \cdot (12k + 2k(k-1)) = 2(12k + 2k^2 - 2k) = 4k^2 + 20k$.
Both also reduce by \vmcompute{}.
\end{proof}

\section{Closed forms for \texorpdfstring{$M_1$ and $M_2$}{M\_1 and M\_2} (Lemma 8.2)}

\begin{definition}[$M_1$ entry]\label{def:M1_entry}
\rocq{PrimeGapS1.MaynardSpec.M1_entry}
\rocqok
\uses{def:G_2, def:K1K2}
For $(b_i, c_i, b_j, c_j)$ basis bidegrees, set $b = b_i + b_j$ and
$c = c_i + c_j$. Then
\[
M_{1, ij} \;=\; \frac{b!}{(K_1 + b + 2c)!} \cdot \Gnk{c}{K_1}.
\]
This is Lemma~8.2 (first part) of arXiv:1311.4600~v3 verbatim, with
$K_1 = k = 105$.
\end{definition}

\begin{definition}[Eq.~8.8 expansion coefficient $\alpha$]\label{def:alpha}
\rocq{PrimeGapS1.MaynardSpec.alpha}
\rocqok
For $b, c, c' \in \N$ with $c' \le c$,
\[
\alpha(b, c, c') \;=\;
   \frac{\binom{c}{c'} \cdot b! \cdot (2c - 2c')!}{(b + 2c - 2c' + 1)!}.
\]
This is the coefficient of
$(1 - \sum_{j \ne 1} t_j)^{b + 2c - 2c' + 1} \cdot (\sum_{j \ne 1} t_j^2)^{c'}$
in the expansion of
$\int_0^{1 - \sum_{j \ne 1} t_j} (1 - P_1)^b P_2^c \, dt_1$,
obtained by expanding $P_2^c = (t_1^2 + \sum_{j \ne 1} t_j^2)^c$
binomially and integrating $t_1^{2(c - c')} (1 - \sum_j t_j)^b$
via the Beta-function identity
$\int_0^X t^a (X - t)^b \, dt = a! b! X^{a + b + 1} / (a + b + 1)!$.
\end{definition}

\begin{definition}[$M_2$ entry]\label{def:M2_entry}
\rocq{PrimeGapS1.MaynardSpec.M2_entry}
\rocqok
\uses{def:alpha, def:G_2, def:K1K2}
For $(b_i, c_i, b_j, c_j)$ basis bidegrees, with
$b'_1 = b_i + 2c_i - 2c'_1 + 1$, $b'_2 = b_j + 2c_j - 2c'_2 + 1$,
$b_{\mathrm{sum}} = b'_1 + b'_2$, $c_{\mathrm{sum}} = c'_1 + c'_2$:
\[
M_{2, ij}^{(1)} \;=\;
   \sum_{c'_1 = 0}^{c_i} \sum_{c'_2 = 0}^{c_j}
     \alpha(b_i, c_i, c'_1)\cdot \alpha(b_j, c_j, c'_2)
     \cdot \frac{b_{\mathrm{sum}}!}{(K_2 + b_{\mathrm{sum}} + 2 c_{\mathrm{sum}})!}
     \cdot \Gnk{c_{\mathrm{sum}}}{K_2}.
\]
This is Lemma~8.2 (second part) of arXiv:1311.4600~v3 for the
\emph{per-coordinate} integrand $J_k^{(1)}$ (not $\sum_m J_k^{(m)}$):
the symmetry factor $k$ is paid in the threshold $4/k$, not in the
matrix entries.
\end{definition}

\begin{remark}
The matrices have up to $36$ rational summands per entry (when
$c_i = c_j = 5$), with intermediate denominators reaching
$\sim 10^{7000}$ digits before the rational normalisation collapses
them. This is why \texttt{MaynardSpec.v} also ships
$\Z$-level twins (\texttt{m1\_num\_den}, \texttt{m2\_num\_den}) that
keep the $(\text{num}, \text{den})$ pair separate; \texttt{MaynardVerify.v}
operates on those twins to avoid \texttt{rat}'s
canonical \texttt{gcd}-normalisation cost during \vmcompute.
\end{remark}

\section{Kernel cross-check: matrices match the spec}

\begin{theorem}[$M_1$ entries match Maynard's closed form]\label{thm:M1_match}
\rocq{PrimeGapS1.MaynardVerify.Def.all_match_M1Z_true}
\rocqok
\uses{def:M1_entry}
For every $(i, j) \in \{0, \ldots, 41\}^2$,
\[
\texttt{m1\_num}(i, j) \cdot \DMone
   \;=\; M_1^{\mathrm{int}}[i][j] \cdot \texttt{m1\_den}(i, j)\quad\text{in }\Z,
\]
where $M_1^{\mathrm{int}}$ is the FLINT-shipped integer matrix and
$(\texttt{m1\_num}, \texttt{m1\_den})$ are the closed-form
$(\text{numerator}, \text{denominator})$ from
Definition~\ref{def:M1_entry}.
\end{theorem}
\begin{proof}\rocqok
The $1764$ Boolean cross-multiplication checks
\texttt{M1\_entry\_matchZ}~$i\,j$ are aggregated into
\texttt{all\_match\_M1Z}~$=$~\texttt{forallb}~over $42 \times 42$, and
the Lemma is closed by a single \vmcompute{}. The check costs
${\sim}90$~s wall-clock. \texttt{Print Assumptions} reports only the
standard \Uintsf{} primitives.
\end{proof}

\begin{theorem}[$M_2$ entries match Maynard's closed form]\label{thm:M2_match}
\rocq{PrimeGapS1.MaynardVerify.all_match_M2Z_true}
\rocqok
\uses{def:M2_entry}
The analogous identity holds for $M_2$:
$\texttt{m2\_num}(i, j) \cdot \DMtwo
   = M_2^{\mathrm{int}}[i][j] \cdot \texttt{m2\_den}(i, j)$
for all $(i, j) \in \{0, \ldots, 41\}^2$.
\end{theorem}
\begin{proof}\rocqok
Same shape as Theorem~\ref{thm:M1_match}, but the entries are double
sums of up to 36 rational terms, so the per-row \vmcompute{} cost is
significant. The check is split into six 7-row chunks
(\texttt{MaynardVerify/M2\_0..5.v}) that \texttt{make -j} runs
concurrently, and the assembly stitches them via \texttt{seq\_split\_42}
plus \texttt{forallb\_app} — no per-entry recomputation in the
\Qed{} of \texttt{all\_match\_M2Z\_true} itself.
\texttt{Print Assumptions}: only the standard \Uintsf{} primitives.
\end{proof}

\begin{remark}
Theorems~\ref{thm:M1_match} and~\ref{thm:M2_match} together close the
trust gap on the input data: every entry of the $42 \times 42$ integer
matrices that the rest of the proof consumes is kernel-verified to
agree with Maynard's closed-form specification (Lemma~8.2). The FLINT
generator is therefore not part of the trust base for the input data
either.
\end{remark}

\section{Paper-form \texorpdfstring{$\leftrightarrow$}{<->} computational-form spec bridge}

The cross-checks above operate on the $\Z$-level twins
(\texttt{m1\_num\_den}, \texttt{m2\_num\_den}) because \vmcompute{} on
$\rat$-level matrices stalls on MathComp's HB canonical-structure
elaboration at concrete dimension $42$. The rat-level spec
(\texttt{M1\_entry}, \texttt{M2\_entry}) is the
documentation-shaped form a reviewer reads side by side with
arXiv:1311.4600~\S8. The bridge below certifies in the kernel that the
two forms encode the same closed forms: a single \Qed{} per matrix,
with no axioms.

\begin{definition}[Reading a $\Z$-pair as a rational]\label{def:qfrac}
\rocq{PrimeGapS1.MaynardSpecBridge.qfrac}
\rocqok
For $(n, d) \in \Z \times \Z$, set
$\texttt{qfrac}(n, d) := \frac{n}{d} \in \rat$
(MathComp's int--to--rat morphism applied component-wise, then divided).
This reads a $(\text{numerator}, \text{denominator})$ pair as the
rational it represents.
\end{definition}

\begin{theorem}[$\Mone$ paper-form spec equals $\Z$-level spec]\label{thm:M1_spec_rat_eq}
\rocq{PrimeGapS1.MaynardSpecBridge.M1_spec_rat_eq}
\rocqok
\uses{def:M1_entry, def:qfrac}
For every $(i, j) \in \{0, \ldots, 41\}^2$,
\[
  \texttt{M1\_spec\_ij}(i, j) \;=\; \texttt{qfrac}\bigl(\texttt{m1\_num\_den\_at}(i, j)\bigr).
\]
\end{theorem}
\begin{proof}\rocqok
Layered structural induction in
\texttt{theories/S1/MaynardSpecBridge.v}: \texttt{factZ\_to\_rat},
\texttt{dblratZ\_to\_rat}, \texttt{binZ\_to\_rat},
\texttt{compositionsZ\_eq\_compositions}, \texttt{G2Z\_to\_rat}, and
\texttt{m1\_num\_den\_to\_rat} compose into a single rewrite chain.
\texttt{Print Assumptions} reports \emph{Closed under the global
context}: no axioms, not even \Uintsf\ primitives.
\end{proof}

\begin{theorem}[$\Mtwo$ paper-form spec equals $\Z$-level spec]\label{thm:M2_spec_rat_eq}
\rocq{PrimeGapS1.MaynardSpecBridge.M2_spec_rat_eq}
\rocqok
\uses{def:M2_entry, def:qfrac}
For every $(i, j) \in \{0, \ldots, 41\}^2$,
\[
  \texttt{M2\_spec\_ij}(i, j) \;=\; \texttt{qfrac}\bigl(\texttt{m2\_num\_den\_at}(i, j)\bigr).
\]
\end{theorem}
\begin{proof}\rocqok
Same layered induction, with two extra steps for the running
$\texttt{qplus}$ accumulator: \texttt{qfrac\_qmul} and
\texttt{qfrac\_qplus} bridge \texttt{qmul} / \texttt{qplus} on
$\Z \times \Z$ pairs to $*$ / $+$ on $\rat$, and a generic
$\texttt{fold\_left qplus} \to \sum$ lemma collects the double sum
over $(c'_1, c'_2)$. \texttt{Print Assumptions}: \emph{Closed under
the global context}.
\end{proof}

\begin{remark}
Theorems~\ref{thm:M1_spec_rat_eq} and~\ref{thm:M2_spec_rat_eq} live
in \texttt{MaynardSpecBridge.v}, a leaf in the dependency DAG:
\texttt{Cert.v} does not import them, and the headline theorem does
not depend on them. Their job is purely to certify that the
documentation-shaped $\rat$-level definitions
(Definitions~\ref{def:M1_entry} and~\ref{def:M2_entry}, written to
match Maynard's paper character for character) and the
$\vmcompute$-shaped $\Z$-level definitions
(\texttt{m1\_num\_den}, \texttt{m2\_num\_den}, consumed by
Theorems~\ref{thm:M1_match} and~\ref{thm:M2_match}) compute the same
rational. Combined with those theorems, this gives a two-step
kernel-verified path from each shipped integer matrix entry back to
Maynard's exact closed form.
\end{remark}

%%%%%%%%%%%%%%%%%%%%%%%%%%%%%%%%%%%%%%%%%%%%%%%%%%%%%%%%%%%%%%%%%%%%%%%%
%%%%%%%%%%%%%%%%%%%%%%%%%%%%%%%%%%%%%%%%%%%%%%%%%%%%%%%%%%%%%%%%%%%%%%%%
\chapter{The Rayleigh-quotient witness route}\label{ch:rayleigh}

This chapter is the analytic heart of the proof. Maynard's
Lemma~8.3 reduces $\Mk[k] > c$ to: exhibit a single
coefficient vector $a \in \rat^{42}$ with
$k \cdot a^T \Mtwo a / a^T \Mone a > c$, equivalently
$c \cdot a^T \Mone a < k \cdot a^T \Mtwo a$ (and $a^T \Mone a > 0$
so the quotient is well-defined). The shipped rational vector
$v_{\mathrm{witness}}$ in \texttt{Witness\_Quad.v} is such an $a$
for $(k, c) = (105, 4)$. The kernel verifies the strict integer
inequality at $v_{\mathrm{witness}}$ by a single \vmcompute{}
reflexivity on $\Z$-arithmetic; the supplementary algebraic shim
lifts that integer fact to the rat-level bigop the headline
surfaces.

The material lives in \texttt{Witness\_Quad.v} (the witness vector
itself) and \texttt{CertRayleigh.v} (the integer Rayleigh inequality,
its positivity companion, and the rat-level lift to the bigop on
the paper-form spec matrices).

\section{The 42-entry rational witness vector}

\begin{definition}[Rational witness $v_{\mathrm{witness}}$]\label{def:v_witness}
\rocq{PrimeGapS1.Witness_Rayleigh.v_witness}
\rocqok
\uses{def:maynard_basis}
$v_{\mathrm{witness}} : \texttt{list }(\Z \times \Z)$ is a list of
42 pairs $(n_i, d_i)$ with $\gcd(|n_i|, d_i) = 1$ and $d_i > 0$,
representing the rational vector $v \in \rat^{42}$ with
$v[i] = n_i / d_i$. The list is autogenerated by
\texttt{python/build\_quad\_witness.py} from the top eigenvector
of $\Mone^{-1}\Mtwo$ (computed via \texttt{acb\_mat.eig} at
$1024$-bit precision), phase-aligned (largest entry positive
real), and snapped to small-denominator rationals via
Mathematica-style absolute-tolerance \texttt{Rationalize} with
$\mathrm{tol} = 10^{-14}$ (continued-fraction convergents). The
entries are listed in the same row/column order as
\texttt{MaynardBasis.maynard\_basis} (Definition~\ref{def:maynard_basis}).
\end{definition}

\begin{remark}[Witness statistics]
The maximum denominator across the 42 entries is
$67\,213\,321\,643\,309 \approx 1.4 \cdot 10^{13}$, i.e.\ about $46$
bits. The eigenvector spans $\approx 14$ decimal orders of magnitude
(smallest non-zero $|v_i| \approx 10^{-14}$, largest $\approx 1$),
which lower-bounds the denominator size: truncating small components
destroys the Rayleigh inequality. The verified slack at this
witness, in exact rationals, is
\[
   \frac{105 \cdot v^T \Mtwo v - 4 \cdot v^T \Mone v}{v^T \Mone v}
   \;\approx\; +2.07 \cdot 10^{-3}.
\]
\end{remark}

\section{Integer reduction of the witness}

\begin{definition}[Common denominator]\label{def:v_den}
\rocq{PrimeGapS1.CertRayleigh.v_den, PrimeGapS1.CertRayleigh.v_num}
\rocqok
\uses{def:v_witness}
$\vden := \mathrm{lcm}\,(d_0, \ldots, d_{41})$ is the LCM of the 42
witness denominators, computed by a left fold over
$v_{\mathrm{witness}}$. The scaled integer numerators
$\vnum : \texttt{list }\Z$ are defined by
$\vnum[i] := n_i \cdot (\vden / d_i)$, so that
$v[i] = \vnum[i] / \vden$ as rationals. This presentation puts
every component over the same denominator, which is what allows the
Rayleigh inequality to reduce to a single integer comparison after
clearing the matrix denominators $\DMone, \DMtwo$.
\end{definition}

\section{Integer Rayleigh inequality}

\begin{definition}[Integer quadratic form]\label{def:quad}
\rocq{PrimeGapS1.CertRayleigh.row_dot, PrimeGapS1.CertRayleigh.mat_vec_mul,
      PrimeGapS1.CertRayleigh.quad}
\rocqok
For $M : \texttt{list (list }\Z\texttt{)}$ and $v : \texttt{list }\Z$,
$\texttt{quad}\,M\,v := v^T \cdot M \cdot v$ (the standard quadratic
form on integer lists), implemented via
$\texttt{row\_dot} : \texttt{list }\Z \to \texttt{list }\Z \to \Z$
(dot product, with the implicit pad-to-shorter convention) and
$\texttt{mat\_vec\_mul}$ (row-wise dot of $M$ with $v$).
\end{definition}

\begin{definition}[Integer numerators of the Rayleigh quotient]\label{def:num_M1M2}
\rocq{PrimeGapS1.CertRayleigh.num_M1, PrimeGapS1.CertRayleigh.num_M2}
\rocqok
\uses{def:quad, def:v_den}
$\texttt{num\_M1} := \texttt{quad}\,M_1^{\mathrm{int}}\,\vnum$ and
$\texttt{num\_M2} := \texttt{quad}\,M_2^{\mathrm{int}}\,\vnum$.
After clearing the common denominator $\DMl \cdot \vden^2$ on both
sides, the strict Rayleigh-quotient bound
$4 \cdot v^T \Mone v < 105 \cdot v^T \Mtwo v$ on the FLINT integer
matrices becomes the strict integer inequality
$4 \cdot \DMtwo \cdot \texttt{num\_M1} < 105 \cdot \DMone \cdot \texttt{num\_M2}$.
\end{definition}

\begin{theorem}[Positivity of the integer numerator]\label{thm:rayleigh_pos}
\rocq{PrimeGapS1.CertRayleigh.rayleigh_witness_M1_positive}
\rocqok
\uses{def:num_M1M2}
$0 < \texttt{num\_M1}$ in $\Z$.
\end{theorem}
\begin{proof}\rocqok
A single \vmcompute{} reflexivity in pure $\Z$-arithmetic.
\texttt{Print Assumptions} reports \emph{Closed under the global
context} — no $\Uintsf$ primitives.
\end{proof}

\begin{theorem}[Strict integer Rayleigh inequality]\label{thm:rayleigh_holds}
\rocq{PrimeGapS1.CertRayleigh.rayleigh_witness_holds}
\rocqok
\uses{def:num_M1M2}
\[
   4 \cdot \DMtwo \cdot \texttt{num\_M1}
      \;<\;
   105 \cdot \DMone \cdot \texttt{num\_M2}
   \quad\text{in } \Z.
\]
\end{theorem}
\begin{proof}\rocqok
A single \vmcompute{} reflexivity on the integer comparison
(${\sim}5$~s on a contemporary machine). Both sides are huge
integers (the matrix denominators reach ${\sim}220$ decimal digits
and the witness denominator $\vden$ has ${\sim}40$ digits), but
the comparison itself is one BigZ subtraction. \texttt{Print
Assumptions} reports \emph{Closed under the global context}.
\end{proof}

\begin{remark}
Theorems~\ref{thm:rayleigh_pos} and~\ref{thm:rayleigh_holds} are
the entire integer-arithmetic content of the headline theorem. If
the Python layer shipped a vector $v_{\mathrm{witness}}$ that failed
the inequality, the kernel \vmcompute{} would fail to reduce to
\texttt{true} and the build would stop. The eigenvector itself is
not in the trust base; only the resulting integer inequality is.
\end{remark}

\section{Lifting to the rat-level Rayleigh bound}

\begin{definition}[Rat-level Rayleigh-quotient numerator]\label{def:quad_spec}
\rocq{PrimeGapS1.CertRayleigh.v_rat, PrimeGapS1.CertRayleigh.quad_spec}
\rocqok
\uses{def:v_den}
For an entry-wise spec $M_{\mathrm{spec}} : \nat \to \nat \to \rat$,
\[
   \texttt{quad\_spec}\,M_{\mathrm{spec}}
   \;:=\;
   \sum_{i < 42} \sum_{j < 42}
      v_{\rat}(i) \cdot M_{\mathrm{spec}}(i,j) \cdot v_{\rat}(j)
   \quad\in\;\rat,
\]
where $v_{\rat}(i) := \texttt{Z2rat}(\vnum[i]) / \texttt{Z2rat}(\vden)$.
This is the standard Rayleigh-quotient numerator $v^T M v$ on the
paper-form spec matrices, expanded as a $\sigb$ bigop (the form the
headline statement consumes).
\end{definition}

\begin{lemma}[Z $\to$ rat bridge for the Rayleigh-quotient sum]
\label{lem:Z2rat_quad}
\rocq{PrimeGapS1.CertRayleigh.Z2rat_quad_eq_sum}
\rocqok
\uses{def:quad, def:v_den}
For any well-formed $M : \texttt{list (list }\Z\texttt{)}$ of
dimension $42 \times 42$ and $v_{\mathrm{num}}$ of length $42$,
\[
   \texttt{Z2rat}\bigl(\texttt{quad}\,M\,\vnum\bigr)
   \;=\;
   \sum_{i < 42} \sum_{j < 42}
      \texttt{Z2rat}(\vnum[i]) \cdot
      \texttt{Z2rat}(\texttt{mat\_get}\,M\,i\,j) \cdot
      \texttt{Z2rat}(\vnum[j]).
\]
\end{lemma}
\begin{proof}\rocqok
Structural induction on the dimension via the
\texttt{row\_dot} / \texttt{mat\_vec\_mul} chain. The base step
$n = 0$ collapses both sides via \texttt{big\_ord0}; the inductive
step combines \texttt{big\_ord\_recl} with the morphism lemmas
$\texttt{Z2rat\_add}$, $\texttt{Z2rat\_mul}$ for the outer
\texttt{row\_dot}, then \texttt{mulr\_sumr} to distribute over the
inner sum.
\end{proof}

\begin{lemma}[Per-cell algebraic identity]\label{lem:quad_cell_identity}
\rocq{PrimeGapS1.CertRayleigh.quad_cell_identity}
\rocqok
For $\texttt{dm}, \texttt{vd}, v_i, m_{ij}, v_j : \rat$ with
$\texttt{vd} \ne 0$ and $\texttt{dm} \ne 0$,
\[
   \texttt{dm} \cdot \texttt{vd}^2 \cdot
     \Bigl(\frac{v_i}{\texttt{vd}} \cdot \frac{m_{ij}}{\texttt{dm}} \cdot
            \frac{v_j}{\texttt{vd}}\Bigr)
   \;=\;
   v_i \cdot m_{ij} \cdot v_j.
\]
\end{lemma}
\begin{proof}\rocqok
A direct \texttt{field} call under the nonzero hypotheses. The lemma
is intentionally \Qed-sealed at top level so that its proof term is
not unfolded inside the $42 \cdot 42 = 1764$ nested binders of the
bigop in Lemma~\ref{lem:quad_spec_eq_Z}.
\end{proof}

\begin{lemma}[Bigop $\leftrightarrow$ Z bridge]\label{lem:quad_spec_eq_Z}
\rocq{PrimeGapS1.CertRayleigh.quad_spec_eq_Z}
\rocqok
\uses{lem:Z2rat_quad, lem:quad_cell_identity,
      thm:M1_match, thm:M2_match,
      thm:M1_spec_rat_eq, thm:M2_spec_rat_eq, lem:M1_spec_eq_int,
      lem:M2_spec_eq_int}
Given a spec $M_{\mathrm{spec}}$ and an integer matrix
$M_{\mathrm{int}}$ such that
$M_{\mathrm{spec}}(i,j) = \texttt{Z2rat}(\texttt{mat\_get}\,M_{\mathrm{int}}\,i\,j) / \texttt{Z2rat}(D_M)$
for all $(i, j) < 42$, with $D_M > 0$ and $M_{\mathrm{int}}$ a
$42 \times 42$ matrix,
\[
   \texttt{Z2rat}(D_M) \cdot \texttt{Z2rat}(\vden)^2 \cdot
      \texttt{quad\_spec}\,M_{\mathrm{spec}}
   \;=\;
   \texttt{Z2rat}(\texttt{quad}\,M_{\mathrm{int}}\,\vnum).
\]
\end{lemma}
\begin{proof}\rocqok
A single application of Lemma~\ref{lem:Z2rat_quad} (to recast the
right-hand side as the rat-level bigop on $\texttt{mat\_get}$ /
$\vnum$) followed by two nested $\texttt{eq\_bigr}$ applications of
Lemma~\ref{lem:quad_cell_identity} to clear the per-cell
denominators $D_M$ and $\vden$. Qed-sealing the per-cell identity
keeps the bigop walk referring to it by name instead of inlining a
fresh \texttt{field} proof term under each of the $42 \cdot 42 =
1764$ binders, so the proof term stays small for kernel
verification. The outer $\texttt{mulr\_sumr}$ rewrite is
restricted to $[LHS]$ to keep unification cheap. \texttt{Print
Assumptions} reports \emph{Closed under the global context}.
\end{proof}

\begin{lemma}[Generic Rayleigh lift]\label{lem:rayleigh_lift_generic}
\rocq{PrimeGapS1.CertRayleigh.rayleigh_lift_generic}
\rocqok
Let $d_1, d_2, a_1, a_2 \in \Z$ with $d_1, d_2 > 0$, $v \in \rat$
with $v > 0$, and $q_1, q_2 \in \rat$. Suppose
\[
\texttt{Z2rat}(d_k) \cdot v^2 \cdot q_k
   \;=\;
   \texttt{Z2rat}(a_k)
\quad\text{for } k = 1, 2,
\]
and the integer comparison
$\texttt{Z2rat}(4 \cdot d_2 \cdot a_1) < \texttt{Z2rat}(105 \cdot d_1 \cdot a_2)$
holds. Then $4 \cdot q_1 < 105 \cdot q_2$ in $\rat$.
\end{lemma}
\begin{proof}\rocqok
Multiply both sides of the goal by the strictly positive scalar
$\texttt{Z2rat}(d_1) \cdot \texttt{Z2rat}(d_2) \cdot v^2$
(monotonicity of $<$ under positive multiplication, $\texttt{ltr\_pM2r}$);
then $\texttt{ring}$ on the two scaled-quad hypotheses rewrites
each side into $\texttt{Z2rat}(4 \cdot d_2 \cdot a_1)$ and
$\texttt{Z2rat}(105 \cdot d_1 \cdot a_2)$ respectively, and the
result is the integer comparison hypothesis. Abstracting over
$q_1, q_2, v$ and the four $\Z$ scalars keeps the two \texttt{ring}
calls operating on tiny abstract scalars rather than on the
concrete \texttt{quad\_spec} bigops, so the proof term stays small.
\texttt{Print Assumptions}: \emph{Closed under the global context}.
\end{proof}

\begin{theorem}[Strict Rayleigh-quotient bound at the witness]
\label{thm:rayleigh_lt_main}
\rocq{PrimeGapS1.CertRayleigh.rayleigh_lt_main}
\rocqok
\uses{lem:quad_spec_eq_Z, lem:rayleigh_lift_generic, thm:rayleigh_holds, thm:rayleigh_pos}
\[
   4 \cdot \texttt{quad\_spec}\,\texttt{M1\_spec\_ij}
   \;<\;
   105 \cdot \texttt{quad\_spec}\,\texttt{M2\_spec\_ij}
   \quad\text{in }\rat.
\]
\end{theorem}
\begin{proof}\rocqok
A one-line \texttt{exact:} application of
Lemma~\ref{lem:rayleigh_lift_generic} to the two specialisations
$\texttt{quad\_M\{1,2\}\_spec\_eq\_Z}$ of
Lemma~\ref{lem:quad_spec_eq_Z} (at $(M_{\mathrm{spec}},
M_{\mathrm{int}}, D_M) = (\texttt{M1\_spec\_ij}, M_1^{\mathrm{int}},
\DMone)$ and the $M_2$ analogue, discharged by
Lemmas~\ref{lem:M1_spec_eq_int} and~\ref{lem:M2_spec_eq_int}),
together with $\texttt{rayleigh\_witness\_holds\_rat}$ — the
integer Rayleigh inequality of Theorem~\ref{thm:rayleigh_holds}
lifted to $\rat$ by $\texttt{Z2rat\_lt}$. \texttt{Print
Assumptions}: \emph{Closed under the global context}.
\end{proof}

\section{The headline}

\begin{theorem}[End-to-end Rayleigh-quotient certification]\label{thm:M105}
\rocq{PrimeGapS1.CertRayleigh.maynard_M105_certified_rayleigh}
\rocqok
\uses{lem:M1_spec_eq_int, lem:M2_spec_eq_int, thm:rayleigh_lt_main}
The conjunction
\[
\begin{aligned}
&\bigl(\forall i, j < 42.\;
  \texttt{M1\_spec\_ij}(i,j)
    = \texttt{Z2rat}(M_1^{\mathrm{int}}[i][j]) \,/\, \texttt{Z2rat}(\DMone)\bigr)\\
&\quad\wedge\;
\bigl(\forall i, j < 42.\;
  \texttt{M2\_spec\_ij}(i,j)
    = \texttt{Z2rat}(M_2^{\mathrm{int}}[i][j]) \,/\, \texttt{Z2rat}(\DMtwo)\bigr)\\
&\quad\wedge\;
   4 \cdot \texttt{quad\_spec}\,\texttt{M1\_spec\_ij}
   \;<\;
   105 \cdot \texttt{quad\_spec}\,\texttt{M2\_spec\_ij},
\end{aligned}
\]
where $\texttt{Z2rat}\,(z : \Z) : \rat := (\texttt{Z\_to\_int}\,z)\%\texttt{:\~{}R}$
is the thin $\Z \to \rat$ embedding.
\end{theorem}
\begin{proof}\rocqok
A three-way \texttt{split}: the two composed per-matrix identities
are discharged by Lemmas~\ref{lem:M1_spec_eq_int}
and~\ref{lem:M2_spec_eq_int} (paper-form spec entry equals FLINT
integer entry over the common denominator), and the third conjunct
is Theorem~\ref{thm:rayleigh_lt_main} (the rat-level Rayleigh
inequality). Modulo Maynard's Lemma~8.3 (paper-side), this entails
$\Mk[105] > 4$. \texttt{Print Assumptions}: \emph{Closed under the
global context}.
\end{proof}

\begin{lemma}[$\Mone$ paper-form spec equals FLINT entry over $\DMone$]\label{lem:M1_spec_eq_int}
\rocq{PrimeGapS1.Cert.M1_spec_eq_int}
\rocqok
\uses{thm:M1_match, thm:M1_spec_rat_eq}
For every $(i, j) \in \{0, \ldots, 41\}^2$,
\[
  \texttt{M1\_spec\_ij}(i, j)
    \;=\; \texttt{Z2rat}\bigl(M_1^{\mathrm{int}}[i][j]\bigr) \,/\,
          \texttt{Z2rat}(\DMone) \quad\text{in } \rat.
\]
\end{lemma}
\begin{proof}\rocqok
Compose Theorem~\ref{thm:M1_spec_rat_eq}, which rewrites the
$\rat$-level spec into $\texttt{qfrac}(\texttt{m1\_num\_den\_at}(i,j))$,
with the $\Z$-level cross-multiplication from
Theorem~\ref{thm:M1_match} (extracted per-entry from the $42 \times 42$
\texttt{forallb}) via the helper
\texttt{qfrac\_eq\_div}: in $\Z$, $a\cdot d = c\cdot b$ with $b, d > 0$
implies $\texttt{qfrac}(a, b) = \texttt{Z2rat}(c)/\texttt{Z2rat}(d)$ in
$\rat$. Denominator positivity is provided by
\texttt{m1\_num\_den\_at\_den\_pos} (product of factorials) and
\texttt{D\_M1\_pos} (\vmcompute).
\end{proof}

\begin{lemma}[$\Mtwo$ paper-form spec equals FLINT entry over $\DMtwo$]\label{lem:M2_spec_eq_int}
\rocq{PrimeGapS1.Cert.M2_spec_eq_int}
\rocqok
\uses{thm:M2_match, thm:M2_spec_rat_eq}
The analogous identity holds for $\Mtwo$:
$\texttt{M2\_spec\_ij}(i, j)
  = \texttt{Z2rat}(M_2^{\mathrm{int}}[i][j]) / \texttt{Z2rat}(\DMtwo)$
for all $(i, j) \in \{0, \ldots, 41\}^2$.
\end{lemma}
\begin{proof}\rocqok
Same shape as Lemma~\ref{lem:M1_spec_eq_int}, with
Theorem~\ref{thm:M2_spec_rat_eq} and Theorem~\ref{thm:M2_match} in
place of their $\Mone$ counterparts.
\end{proof}

\paragraph{Reading.}
A reviewer who trusts \emph{only Rocq's kernel} obtains, from
Theorem~\ref{thm:M105}, a certified strict Rayleigh-quotient bound
on the paper-form spec matrices at the shipped 42-entry rational
witness vector, together with the closed-form match of those
matrices to Maynard's specification. A reviewer who additionally
accepts Maynard's Lemma~8.3 (an Annals-published Rayleigh-quotient
identity, paper-side) obtains $\Mk[105] > 4$ end-to-end.

%%%%%%%%%%%%%%%%%%%%%%%%%%%%%%%%%%%%%%%%%%%%%%%%%%%%%%%%%%%%%%%%%%%%%%%%
%%%%%%%%%%%%%%%%%%%%%%%%%%%%%%%%%%%%%%%%%%%%%%%%%%%%%%%%%%%%%%%%%%%%%%%%
\chapter{The eigenvalue route (target form)}\label{ch:eigen}

The Rayleigh chapter certifies a strict inequality at one witness vector;
combined with Maynard's Lemma~8.3 (paper-side) that already gives
$\Mk[105] > 4$. This chapter establishes a \emph{second}, self-contained
Rocq result that matches Maynard's own \emph{characterisation} of $\Mk$ as
($k$ times) the largest eigenvalue of the generalised eigenproblem
$\Mtwo v = \lambda \Mone v$: it exhibits, inside the kernel, an eigenvalue
$> 4$ of the rational matrix
$M_{105} = 105 \cdot \Mone^{-1}\Mtwo$. The headline of this route is
\rocq{PrimeGapS1.MaynardEigen.maynard_M105_certified}
(Theorem~\ref{thm:M105_eigen}).

This route is built \emph{on} the Rayleigh witness — the strict inequality
of Theorem~\ref{thm:rayleigh_lt_main} is exactly the positivity input the
spectral bridge consumes — so Route~1 is at once a direct proof of the
inequality and the computational core of Route~2. Everything below is
\Qed{} and axiom-free; in particular this route genuinely \emph{does}
build a characteristic polynomial (modulo primes), perform a
Chinese-remainder lift, and exhibit an eigenvalue, all over Stdlib's $\Z$
and $\Fp$ with no \texttt{native\_compute} and no real-algebra
(IVT / Sturm / $\realalg$) machinery.

\section{The generalised-eigenvalue matrices}

\begin{definition}[Rational pencil matrices]\label{def:eigen_mats}
\rocq{PrimeGapS1.EigenBridge.M1_rat, PrimeGapS1.EigenBridge.M2_rat,
      PrimeGapS1.EigenBridge.A_rat, PrimeGapS1.EigenBridge.M105}
\rocqok
As genuine MathComp matrices in $\mathrm{'M}[\rat]_{42}$:
$\Mone^{\rat} := M_1^{\mathrm{int}} / \DMone$ and
$\Mtwo^{\rat} := M_2^{\mathrm{int}} / \DMtwo$ (entrywise int-to-rat
division), with $A^{\rat} := (\Mone^{\rat})^{-1} \Mtwo^{\rat}$ and
$M_{105} := 105 \cdot A^{\rat}
   = 105 \cdot (\Mone^{\rat})^{-1}\Mtwo^{\rat}$.
\end{definition}

\begin{definition}[Closed-form trust contract]\label{def:matches_closed_forms}
\rocq{PrimeGapS1.EigenBridge.matches_closed_forms}
\rocqok
\uses{def:M1_entry, def:M2_entry, def:eigen_mats}
$\texttt{matches\_closed\_forms}(M)$ holds when $M = M_{105}$ and the
paper-form specs match the FLINT integer data entrywise:
$\texttt{M1\_spec\_ij}(i,j)
   = \texttt{Z2rat}(M_1^{\mathrm{int}}[i][j])/\texttt{Z2rat}(\DMone)$
and the $\Mtwo$ analogue, for all $(i,j) < 42$.
\end{definition}

\begin{lemma}[The pencil meets the contract]\label{lem:matches_closed_forms_M105}
\rocq{PrimeGapS1.EigenBridge.matches_closed_forms_M105}
\rocqok
\uses{def:matches_closed_forms, lem:M1_spec_eq_int, lem:M2_spec_eq_int}
$\texttt{matches\_closed\_forms}(M_{105})$.
\end{lemma}
\begin{proof}\rocqok
The first conjunct is reflexivity; the other two are
Lemmas~\ref{lem:M1_spec_eq_int} and~\ref{lem:M2_spec_eq_int}.
\end{proof}

\section{The Hermitian variational crux}

The bridge from the Rayleigh inequality to a genuine eigenvalue runs
through a single spectral fact, proved over a \texttt{numClosedFieldType}
(instantiated at $\algC$, MathComp's field of algebraic complex numbers)
on top of the \emph{complex} spectral theorem shipped in MathComp~2.5.0
(\texttt{orthomx\_spectralP}). No real spectral theorem (mathcomp
PR~\#1611) and no project axiom is used.

\begin{theorem}[Hermitian crux]\label{thm:herm_crux}
\rocq{PrimeGapS1.SpectralCrux.herm_crux}
\rocqok
For $A \in \mathrm{'M}[\algC]_n$ Hermitian
($A \in \texttt{hermsymmx}$) and a vector
$w \in \mathrm{'rV}[\algC]_n$ with $(w\,A\,w^{*})_{00} > 0$, there exists
$a$ with $\texttt{eigenvalue}\,A\,a$ and $a > 0$: a Hermitian matrix whose
Hermitian quadratic form is strictly positive at some vector has a
strictly positive eigenvalue.
\end{theorem}
\begin{proof}\rocqok
Spectral-decompose $A = P^{-1} (\mathrm{diag}\,\mathrm{sp})\,P$ with $P$
unitary and the spectrum $\mathrm{sp}$ real
(\texttt{hermitian\_spectral\_diag\_real}). The form rewrites as
$\sum_j \mathrm{sp}_{0j}\,|u_j|^2$ with $u := P\,w^{*}$; strict positivity
forces some $\mathrm{sp}_{0j} > 0$ (\texttt{realsum\_pos}), and that
diagonal entry is an eigenvalue (\texttt{eig\_transfer}). Only
$\texttt{hermsymmx}$ (not $\texttt{realmx}$) is required, so a
\emph{complex} congruence factor may be used downstream — no real Cholesky
is needed.
\end{proof}

\section{Axiom-free positive definiteness of \texorpdfstring{$\Mone$}{M\_1}}

The crux needs a Hermitian \emph{congruence factor} of $\Mone^{\rat}$ over
$\algC$, i.e.\ $\Mone^{\rat}$ positive definite. The certificate is a
characteristic-polynomial sign argument carried out entirely in $\Z$ and
$\Fp$, reconstructed by a deterministic Chinese remainder theorem — not a
floating-point or probabilistic check.

\subsection{The list-\texorpdfstring{$\Z$}{Z} characteristic polynomial}

\begin{definition}[Integer characteristic polynomial]\label{def:char_poly_int}
\rocq{PrimeGapS1.CharPoly.char_poly_int}
\rocqok
$\texttt{char\_poly\_int}(M) : \texttt{list }\Z$ is the coefficient list of
$\charpolysym_M(x) = \det(xI - M)$ for an integer matrix $M$, built on the
list-$\Z$ polynomial library \texttt{IntPoly.v} (which is therefore part of
the eigenvalue route, not legacy). For the $42 \times 42$ matrix
$\Mone^{\mathrm{int}}$ it has $43$ coefficients.
\end{definition}

\subsection{Modular char-poly: Faddeev--LeVerrier and fast Hessenberg}

\begin{theorem}[Modular Faddeev--LeVerrier is sound]\label{thm:char_poly_modZ_sound}
\rocq{PrimeGapS1.ModularFL.char_poly_modZ, PrimeGapS1.ModularFL.char_poly_modZ_sound}
\rocqok
\uses{def:char_poly_int}
For a prime $p$, a square matrix $M$ with $\#M + 1 < p$ and the
Faddeev--LeVerrier divisibility side-condition,
$\texttt{char\_poly\_modZ}\,p\,M
   = \texttt{map}\,(\cdot \bmod p)\,(\texttt{char\_poly\_int}\,M)$.
\end{theorem}

\begin{theorem}[Fast Hessenberg char-poly is sound]\label{thm:char_poly_hess_sound}
\rocq{PrimeGapS1.ModularHess.char_poly_hess, PrimeGapS1.ModularHess.char_poly_hess_sound}
\rocqok
\uses{def:char_poly_int, thm:char_poly_modZ_sound}
Under the same hypotheses,
$\texttt{char\_poly\_hess}\,p\,M
   = \texttt{map}\,(\cdot \bmod p)\,(\texttt{char\_poly\_int}\,M)$.
The $O(n^3)$ Hessenberg pass replaces the $O(n^4)$ Faddeev--LeVerrier pass
for the heavy per-prime work; it returns \texttt{None} on a vanishing pivot
and then falls back to the (sound) FL pass, so no good-pivot hypothesis is
needed.
\end{theorem}
\begin{proof}\rocqok
The \texttt{None} branch reuses Theorem~\ref{thm:char_poly_modZ_sound}.
The \texttt{Some} branch factors through two structural lemmas, both
\emph{fully proven and axiom-free}:
\rocq{PrimeGapS1.ModularHess.hess_recurrence_sound} (the upper-Hessenberg
leading-principal-minor determinant recurrence computes $\charpolysym$, by
Laplace expansion along the last row) and
\rocq{PrimeGapS1.ModularHess.hess_reduce_similar} ($\charpolysym$ is
invariant under the elementary conjugations of the Hessenberg reduction,
i.e.\ that reduction is a similarity), transported coefficientwise through
the ring homomorphism $\Z \to \Fp$.
\end{proof}

\subsection{The CRT reconstruction}

\begin{lemma}[A-priori coefficient bound]\label{lem:char_poly_int_coeff_bound}
\rocq{PrimeGapS1.Bound.char_poly_int_coeff_bound}
\rocqok
\uses{def:char_poly_int}
Every coefficient of $\texttt{char\_poly\_int}(M)$ is bounded in absolute
value by a Hadamard-style quantity
$\texttt{fl\_coeff\_bound}(\#M)(\texttt{max\_abs\_entry}\,M)$ that depends
only on the dimension and the largest entry magnitude.
\end{lemma}

\begin{theorem}[Deterministic CRT reconstruction]\label{thm:crt_reconstruct}
\rocq{PrimeGapS1.CRTCheck.crt_reconstruct}
\rocqok
For integers $c, d$ and a list $ps$ of distinct primes: if every $p \in ps$
divides $c - d$ and $2\,|c - d| < \prod ps$, then $c = d$.
\end{theorem}

\begin{lemma}[Per-prime agreement]\label{lem:per_prime}
\rocq{PrimeGapS1.CRTFrame.per_prime_hess_all, PrimeGapS1.CRTFrame.per_prime_ok}
\rocqok
\uses{thm:char_poly_hess_sound}
For every prime $p$ of the table \texttt{crt\_primes\_M1},
$\texttt{char\_poly\_hess}\,p\,M_1^{\mathrm{int}}
   = \texttt{map}\,(\cdot \bmod p)\,\texttt{cp\_M1\_value}$
(\texttt{per\_prime\_hess\_all}), and the same equality holds for the
Faddeev--LeVerrier pass (\texttt{per\_prime\_ok}).
\end{lemma}
\begin{proof}\rocqok
The aggregate Boolean \texttt{forallb} over the $200$ table primes
(\texttt{WitnessM1CharPoly.crt\_primes\_M1}, each just under $2^{30}$) is
closed by \vmcompute, sharded across eight files
(\texttt{CRTFrame\_part0..7.v}) so \texttt{make -j8} runs the heavy
$O(n^3)$ reductions concurrently and recombines them via
\texttt{forallb\_app} with no per-prime recomputation in the
\Qed. \texttt{per\_prime\_ok} transports the fast-path equality to the FL
pass through Theorems~\ref{thm:char_poly_hess_sound}
and~\ref{thm:char_poly_modZ_sound}.
\end{proof}

\begin{theorem}[$\charpolysym_{\Mone}$ and its sign pattern]\label{thm:cp_M1}
\rocq{PrimeGapS1.M1CharPoly.char_poly_int_M1_eq, PrimeGapS1.M1CharPoly.cp_M1_alternates}
\rocqok
\uses{thm:char_poly_modZ_sound, thm:crt_reconstruct,
      lem:char_poly_int_coeff_bound, lem:per_prime}
$\texttt{char\_poly\_int}(M_1^{\mathrm{int}}) = \texttt{cp\_M1\_value}$
(the shipped $43$-coefficient list), and \texttt{cp\_M1\_value} has
strictly alternating signs with positive constant and leading coefficients.
\end{theorem}
\begin{proof}\rocqok
Coefficient by coefficient: for each index and each table prime the two
$k$-th coefficients agree modulo $p$ (Lemma~\ref{lem:per_prime} with
Theorem~\ref{thm:char_poly_modZ_sound}); the a-priori bound
(Lemma~\ref{lem:char_poly_int_coeff_bound}) lies below the prime product
(a single \vmcompute), so Theorem~\ref{thm:crt_reconstruct} collapses the
$200$ congruences into an honest equality over $\Z$. The sign pattern is
then a \vmcompute{} on the concrete list.
\end{proof}

\subsection{From sign alternation to a congruence factor}

\begin{theorem}[$\Mone$ is positive definite]\label{thm:M1_rat_factor}
\rocq{PrimeGapS1.M1PosDef.M1_rat_factor}
\rocqok
\uses{thm:cp_M1, thm:herm_crux}
There is a unitary-scaled congruence factor
$R \in \mathrm{'M}[\algC]_{42}$ with $R \in \texttt{unitmx}$ and
$R^{*} R = \texttt{map\_mx ratr}\,\Mone^{\rat}$.
\end{theorem}
\begin{proof}\rocqok
A monic polynomial over $\algC$ with integer, strictly sign-alternating
coefficients (positive constant) is strictly positive at every real
$x \le 0$, so each of its real roots is $> 0$. By the complex spectral
theorem the spectrum of the Hermitian $\texttt{map\_mx ratr}\,\Mone^{\rat}$
is real and consists of such roots (Theorem~\ref{thm:cp_M1}), hence
positive; the diagonal matrix of the $\sqrt{\cdot}$ of the spectrum, times
the spectral basis, is the factor $R$, after transferring the $1/\DMone$
scaling over $\rat$ (no division in $\algC$).
\end{proof}

\section{The spectral bridge and the headline}

\begin{lemma}[Abstract spectral bridge]\label{lem:eigen_bridge}
\rocq{PrimeGapS1.MaynardEigen.eigen_bridge}
\rocqok
\uses{thm:herm_crux}
Let $M_1^c, M_2^c, R, S', A \in \mathrm{'M}[\algC]_n$ and
$u \in \mathrm{'cV}[\algC]_n$ with $R \in \texttt{unitmx}$,
$R^{*} R = M_1^c$, $M_2^c$ Hermitian,
$S' = 105 \cdot M_2^c - 4 \cdot M_1^c$,
$A = 105 \cdot (M_1^c)^{-1} M_2^c$, and $(u^{*} S'\, u)_{00} > 0$. Then $A$
has an eigenvalue $> 4$.
\end{lemma}
\begin{proof}\rocqok
$S_2 := (R^{-1})^{*} S'\, R^{-1}$ is Hermitian (congruence) and the witness
$w := (R\,u)^{*}$ makes its Hermitian form positive; Theorem~\ref{thm:herm_crux}
yields a positive eigenvalue $a$ of $S_2$. Since $S_2$ is similar to
$(M_1^c)^{-1} S' = A - 4\cdot I$, the shifted value $a + 4$ is an
eigenvalue of $A$, with $a + 4 > 4$. The lemma is stated on abstract
matrices, so its kernel type-check never forces the
$\sim 10^{200}$-denominator FLINT data.
\end{proof}

\begin{theorem}[End-to-end eigenvalue certification (target form)]\label{thm:M105_eigen}
\rocq{PrimeGapS1.MaynardEigen.maynard_M105_certified}
\rocqok
\uses{lem:matches_closed_forms_M105, thm:M1_rat_factor, lem:eigen_bridge,
      thm:rayleigh_lt_main}
With $\texttt{ratrM} := \texttt{map\_mx}\,(\texttt{ratr} : \rat \to \algC)$
the entrywise embedding of rational matrices into $\algC$:
\begin{verbatim}
Theorem maynard_M105_certified :
  matches_closed_forms M105 /\
  exists lam : algC, eigenvalue (ratrM M105) lam /\ (4 < lam).
\end{verbatim}
i.e.\ the closed-form contract holds and
$M_{105} = 105 \cdot \Mone^{-1}\Mtwo$ has an eigenvalue $> 4$ over $\algC$ —
Maynard's characterisation of $\Mk[105]$ as ($105$ times) the largest
eigenvalue of the generalised eigenproblem $\Mtwo v = \lambda \Mone v$.
\end{theorem}
\begin{proof}\rocqok
The first conjunct is Lemma~\ref{lem:matches_closed_forms_M105}. For the
second, the complex congruence factor of Theorem~\ref{thm:M1_rat_factor}
and the Rayleigh positivity of Theorem~\ref{thm:rayleigh_lt_main} (pushed
through \texttt{ratr}, where it becomes the positive Hermitian form
feeding the crux) discharge the hypotheses of the bridge
Lemma~\ref{lem:eigen_bridge}. \texttt{Print Assumptions}: \emph{Closed
under the global context}.
\end{proof}

\paragraph{Reading.}
Route~2 is the \emph{target form} of the result: rather than a strict
Rayleigh inequality at one witness (Route~1), it states the conclusion in
Maynard's own spectral language — $M_{105}$ has an eigenvalue exceeding
$4$. The two routes share the same axiom-free, $\Z$/$\Fp$-only
computational footprint, and Route~2's only numerical input beyond the
matrix data is the very inequality Route~1 certifies.

%%%%%%%%%%%%%%%%%%%%%%%%%%%%%%%%%%%%%%%%%%%%%%%%%%%%%%%%%%%%%%%%%%%%%%%%
\chapter{Trust base}\label{ch:trust}

\section{What is a Qed and what is paper-side}

\begin{itemize}
\item Every claim of the form $a = b$ in $\Z$ or $\Q$, $a < b$ in
  $\Z$ or $\Q$, or membership of a rational in an interval, that
  appears in the dependency tree of Theorem~\ref{thm:M105} is
  \Qed{} in the source tree. There are no \texttt{Admitted}, no
  \texttt{Axiom}, and no \texttt{Parameter} declarations anywhere
  in \texttt{theories/S1/}.
\item The matrix entries themselves (Lemma~8.2 of Maynard) are
  \Qed{} via Theorems~\ref{thm:M1_match} and~\ref{thm:M2_match}:
  the $1764 + 1764$ kernel-checked cross-multiplications close the
  trust loop on the input data.
\item Maynard's Lemma~8.3 ($\Mk = k \cdot \sup_F J_k(F)/I_k(F)$) is
  the generalised Rayleigh-quotient identity that bridges the
  strict bound at any individual $a \in \rat^{42}$ to $\Mk > 4$.
  This is the analytic step proved in the Annals paper (refereed
  there). It is \emph{not} formalised in this repository, by design
  of the project's scope.
\item The analytic Beta-integral derivation that produces Maynard's
  closed-form $M_1$ and $M_2$ entries (Lemma~8.1 / 8.2 in the paper)
  is taken as a \emph{definition} in Rocq. The kernel certifies that
  the shipped integer matrices match those closed forms, not that
  the closed forms themselves are correct integrations.
\end{itemize}

\section{Reported assumptions}

\emph{Both} headline theorems report a single line under
\texttt{Print Assumptions} — \emph{Closed under the global context}:
\rocq{PrimeGapS1.CertRayleigh.maynard_M105_certified_rayleigh} (Route~1,
the direct Rayleigh witness, Theorem~\ref{thm:M105}) and
\rocq{PrimeGapS1.MaynardEigen.maynard_M105_certified} (Route~2, the
eigenvalue / target form, Theorem~\ref{thm:M105_eigen}). This is stronger
than the typical \vmcompute{} footprint: because the matrices and the
witness vector are encoded as \texttt{list (list Z)} and
\texttt{list (Z * Z)} rather than as native 63-bit packed arrays,
\vmcompute{} reduces through $\Z$ arithmetic — and, for the CRT
positive-definiteness certificate of Route~2, through $\Fp = \Z/p\Z$ —
instead of through the $\Uintsf$ kernel primitives. None of
\texttt{PrimInt63}, \texttt{Uint63Axioms}, or \texttt{CarryType} appears
in the reported assumptions of either theorem.

The project does not invoke \texttt{native\_compute} anywhere — a
\texttt{grep -n native\_compute theories/S1} returns zero hits in
real proof code.

\section{The FLINT layer is outside the trust base}

The FLINT pipeline (Python + \texttt{python-flint}) is the
candidate generator for the matrix entries and the Rayleigh-
quotient witness, plus an independent cross-check, but the Rocq
proof never invokes Python and never loads the JSON certificate
directly. The autogenerated \texttt{Witness.v} and
\texttt{Witness\_Quad.v} are ordinary Rocq files; if the FLINT
layer shipped wrong data, one of the following \vmcompute{}-based
checks would fail:
\begin{itemize}
\item Theorems~\ref{thm:M1_match} and~\ref{thm:M2_match}: the
  $42 \times 42$ matrix entries match Maynard's closed form
  (Lemma~8.2).
\item Theorem~\ref{thm:rayleigh_pos}: positivity of the integer
  Rayleigh numerator $v^T M_1^{\mathrm{int}} v > 0$ at the shipped
  witness.
\item Theorem~\ref{thm:rayleigh_holds}: the strict integer
  Rayleigh inequality at the shipped witness.
\end{itemize}
Each is a \Qed{} closed by pure kernel arithmetic.

%%%%%%%%%%%%%%%%%%%%%%%%%%%%%%%%%%%%%%%%%%%%%%%%%%%%%%%%%%%%%%%%%%%%%%%%
\chapter{Numerical highlights}\label{ch:numeric}

\paragraph{Matrix dimension.} $42 \times 42$.
Maynard's $\Mk[105]$ construction uses $42$ basis polynomials
$\{F_{b,c} : b + 2c \le 11\}$.

\paragraph{Denominators.} $\DMone$ has ${\sim}221$ decimal digits,
$\DMtwo$ ${\sim}227$. Both are strictly positive
(\vmcompute{}).

\paragraph{Witness denominators.}
The $42$ rational entries of $v_{\mathrm{witness}}$
(Definition~\ref{def:v_witness}) have denominators bounded by
$67\,213\,321\,643\,309 \approx 1.4 \cdot 10^{13}$ (${\sim}14$
decimal digits, $46$ bits). The smallest non-zero $|v_i|$ is
$\approx 10^{-14}$ and the largest is $\approx 1$, so the
eigenvector spans ${\sim}14$ decimal orders of magnitude — this
lower-bounds the denominator size, since truncating small
components destroys the inequality.

\paragraph{Verified slack.}
$(105 \cdot v^T \Mtwo v - 4 \cdot v^T \Mone v) / v^T \Mone v
   \approx +2.07 \cdot 10^{-3}$.

\paragraph{Build profile.}
A clean \texttt{make -j8} rebuild takes ${\sim}8.6$~min wall-clock and
produces $43$ \texttt{.vo} files. It splits into two largely concurrent
parts:
\begin{itemize}
\item \emph{Matrix / spec / Rayleigh base} (${\sim}4.5$~min after a
  \texttt{vm\_cast\_no\_check} speed-up), dominated by the reflection
  mass-checks. Theorem~\ref{thm:M2_match}
  (\texttt{all\_match\_M2Z\_true}) — $1764$ entries, each up to $36$
  rational summands with intermediate denominators reaching
  $\sim 10^{7000}$ digits — is split into six $7$-row chunks
  (\texttt{MaynardVerify/M2\_0..5.v}) that \texttt{make -j} runs in
  parallel and stitches via \texttt{forallb\_app}. Theorem~\ref{thm:M1_match}
  and Theorem~\ref{thm:rayleigh_holds} are comparatively cheap, and
  \texttt{Witness.v} parsing is a one-off \BigZ\ literal cost.
\item \emph{Eigenvalue chain} (${\sim}2.3$~min): the $200$-prime CRT
  positive-definiteness certificate for $\Mone$ (Lemma~\ref{lem:per_prime}),
  whose per-prime $O(n^3)$ Hessenberg \vmcompute{} is sharded across
  \texttt{CRTFrame\_part0..7.v} for \texttt{make -j8} and recombined via
  \texttt{forallb\_app} (serially this step is ${\sim}36$~min).
\end{itemize}

\paragraph{Project size.} $43$ \texttt{.v} files under
\texttt{theories/S1/} ($36$ top-level plus $7$ in the
\texttt{MaynardVerify/} parallel-chunk subdirectory), ${\sim}15{,}300$
lines total; \texttt{Witness.v} alone is ${\sim}5{,}700$ lines of
autogenerated certificate data, \texttt{WitnessM1CharPoly.v} ships the
$200$ CRT primes and the $43$ char-poly coefficients, and
\texttt{Witness\_Rayleigh.v} is ${\sim}80$ lines of $42$
$(\mathrm{num}, \mathrm{den})$ pairs.

%%%%%%%%%%%%%%%%%%%%%%%%%%%%%%%%%%%%%%%%%%%%%%%%%%%%%%%%%%%%%%%%%%%%%%%%
\chapter{Map of key lemmas and files}\label{ch:map}

The dependency graph is essentially linear, with one branch:

\paragraph{Matrix-spec backbone.}
\texttt{Witness} $\to$ \texttt{MaynardFactQ} $\to$
\texttt{MaynardBasis} + \texttt{MaynardSpec} $\to$
\texttt{MaynardVerify/} (\texttt{Def.v} + \texttt{M2\_0..5.v} +
assembly \texttt{MaynardVerify.v}) +
\texttt{MaynardSpecBridge.v}.
Headline outputs: Theorems~\ref{thm:M1_match}, \ref{thm:M2_match},
\ref{thm:M1_spec_rat_eq}, \ref{thm:M2_spec_rat_eq}, composed by
\texttt{Cert.v} (Lemmas~\ref{lem:M1_spec_eq_int},
\ref{lem:M2_spec_eq_int}).

\paragraph{Rayleigh-witness backbone.}
\texttt{Witness\_Rayleigh} $\to$ \texttt{CertRayleigh}.
Headline outputs: Theorems~\ref{thm:rayleigh_pos},
\ref{thm:rayleigh_holds}, Lemmas~\ref{lem:quad_spec_eq_Z},
\ref{lem:rayleigh_lift_generic}, Theorem~\ref{thm:rayleigh_lt_main}.

\paragraph{Eigenvalue-route backbone.}
\texttt{SpectralCrux} (the Hermitian crux) together with the $\Z$/$\Fp$
char-poly stack — \texttt{IntPoly} $\to$ \texttt{CharPoly},
\texttt{ModularFL}, \texttt{ModularHess}, \texttt{Bound}, \texttt{CRTCheck},
\texttt{WitnessM1CharPoly}, \texttt{CRTFrameDefs} +
\texttt{CRTFrame\_part0..7} $\to$ \texttt{CRTFrame} $\to$
\texttt{M1CharPoly} $\to$ \texttt{M1PosDef} — feed \texttt{EigenBridge} and
\texttt{MaynardEigen}. Headline output: Theorem~\ref{thm:M105_eigen}
(\texttt{maynard\_M105\_certified}), which also consumes the Rayleigh
inequality Theorem~\ref{thm:rayleigh_lt_main}.

The matrix-spec and Rayleigh-witness backbones merge at
\texttt{CertRayleigh.v}, producing the Route~1 headline
Theorem~\ref{thm:M105} (\texttt{Require Import PrimeGapS1.CertRayleigh.}
loads it); the eigenvalue-route backbone builds on that inequality to
produce the Route~2 headline Theorem~\ref{thm:M105_eigen}
(\texttt{Require Import PrimeGapS1.MaynardEigen.} loads the whole proof).
